\documentclass[11pt]{article}
\usepackage{amsmath, amssymb, amsthm}
\usepackage{geometry}
\usepackage{booktabs}
\usepackage{float}
\floatstyle{ruled}
\newfloat{algorithm}{H}{loa}
\floatname{algorithm}{Algorithm}
\usepackage{hyperref}
\geometry{margin=1in}

\newcommand{\tv}{\tilde{v}}
\newcommand{\te}{\tilde{e}}
\newcommand{\bsbeta}{\boldsymbol{\beta}}

\title{Finding the Optimal Tau Function via Grid Search}
\date{}

\begin{document}
\maketitle

\section{Overview}

The goal is to find a polynomial threshold function
$\tau : \mathbb{R} \to \mathbb{R}$
such that running the collateralized auction with threshold $\tau$
maximizes expected social welfare relative to the VCG benchmark.
An advertiser with value--externality pair $(v, e)$
is admitted to the collateralized auction only if $v \geq \tau(e)$.

We search over the coefficients of $\tau$ exhaustively on a finite grid,
evaluating every candidate against a shared set of simulated auction draws.
The entire evaluation is performed in a normalized coordinate system
in which both dimensions span $[-1, 1]$,
and the best coefficients are then converted back to original units.

\section{Data Preparation}

\subsection{Cost Scaling}

Starting from the raw scores $v_t^{\text{score}}$ (actions per month) and
$e_t^{\text{score}}$ (normalized externality per month) described in
\texttt{value\_externality\_definitions.tex}, the advertiser value and
externality are obtained by applying cost parameters:
\begin{equation}
    v_t = \zeta_v \cdot v_t^{\text{score}}, \qquad
    e_t = \zeta_e \cdot e_t^{\text{score}},
    \label{eq:cost_scaling}
\end{equation}
where $\zeta_v$ (\texttt{action\_cost}, default 1.0) and
$\zeta_e$ (\texttt{externality\_cost}) are free parameters.

\subsection{Independent Normalization}

Each column is mapped to $[-1, 1]$ independently using its own minimum and maximum.
Let $v_{\min}, v_{\max}$ and $e_{\min}, e_{\max}$ denote the column-wise extrema
over the full dataset. Define the range of each column:
\begin{equation}
    R_v = v_{\max} - v_{\min}, \qquad R_e = e_{\max} - e_{\min}.
\end{equation}
The normalized values are:
\begin{equation}
    \tv_t = \frac{2(v_t - v_{\min})}{R_v} - 1 \in [-1,1], \qquad
    \te_t = \frac{2(e_t - e_{\min})}{R_e} - 1 \in [-1,1].
    \label{eq:normalization}
\end{equation}
Joint normalization (a single shared min/max across both columns) is
deliberately avoided: when $\zeta_e \ll \zeta_v$, the $e$ values are far
smaller in magnitude than $v$, so a joint scale would compress $\te$ to
near zero and leave the search unable to exploit the externality dimension.
Independent normalization ensures both dimensions always span the full $[-1,1]$
range regardless of the cost parameters.

\subsection{Welfare Objective Equivalence}

Because $v$ and $e$ are normalized with different scales, a fitness function
that sums $\tv + \te$ would give each dimension equal weight in normalized
space, which does not reflect their relative magnitudes in original units.
To recover the correct relative weighting, we use a scaled objective.

The original welfare of a set $\mathcal{W}$ of winning advertisers is
\begin{equation}
    W(\mathcal{W}) = \sum_{i \in \mathcal{W}} v_i + \sum_{i \in \mathcal{W}} e_i.
    \label{eq:original_welfare}
\end{equation}
Substituting the inverse of \eqref{eq:normalization},
$v_i = \tv_i \cdot R_v/2 + (v_{\min}+v_{\max})/2$
and analogously for $e_i$, and dropping additive constants
that do not affect the argmax, the equivalent normalized objective is
\begin{equation}
    F_{\text{norm}}(\mathcal{W})
    = \sum_{i \in \mathcal{W}} \tv_i
    + \underbrace{\frac{R_e}{R_v}}_{\lambda} \sum_{i \in \mathcal{W}} \te_i.
    \label{eq:scaled_welfare}
\end{equation}
The scale factor $\lambda = R_e/R_v$ preserves the relative importance of
the two dimensions. Since $R_e = \zeta_e \cdot (e_{\max}^{\text{score}} - e_{\min}^{\text{score}})$
and $R_v = \zeta_v \cdot (v_{\max}^{\text{score}} - v_{\min}^{\text{score}})$,
$\lambda$ encodes the externality cost parameter $\zeta_e$ together with
the natural scale difference between action and externality scores.
Maximizing \eqref{eq:scaled_welfare} in normalized space is equivalent to
maximizing \eqref{eq:original_welfare} in original space.

\section{Advertiser Sampling}

For each of $K$ simulated auctions, $n$ tweets are drawn without replacement
from the full dataset using a fixed random seed. Each sampled tweet $i$
yields two advertiser pairs: the original-space pair $(v_i, e_i)$ and the
normalized pair $(\tv_i, \te_i)$. Denote the normalized auction draws as
$\mathcal{A} = \bigl\{ A^{(j)} \bigr\}_{j=1}^{K}$,
where each $A^{(j)} = \bigl\{(\tv_i^{(j)}, \te_i^{(j)})\bigr\}_{i=1}^{n}$.
The search is conducted entirely with the normalized pairs.

\section{Grid Construction}

We search over degree-$d$ polynomials of the form
\begin{equation}
    \tau(\te;\,\bsbeta) = \sum_{p=0}^{d} \beta_p\, \te^{\,p},
    \qquad \bsbeta = (\beta_0, \ldots, \beta_d) \in \mathbb{R}^{d+1}.
\end{equation}
A discrete grid is constructed by taking the Cartesian product of
one-dimensional ranges for each coefficient:
\begin{align}
    \beta_0 &\in \mathcal{B}_0 = \text{linspace}(-2,\; 2,\; M), \\
    \beta_p &\in \mathcal{B} = \text{linspace}(-5,\; 5,\; M), \quad p = 1, \ldots, d,
\end{align}
where $M$ is the number of grid points per dimension (\texttt{num\_points}, default 20).
The intercept range $[-2, 2]$ covers the scale of $\tilde{v} \in [-1,1]$.
The slope range $[-5, 5]$ covers meaningful separations: with $\te \in [-1,1]$,
a coefficient $|\beta_p| = 5$ shifts $\tau$ by at most 5 across the full
externality range, which is large relative to the data scale.

The full grid is the Cartesian product
\begin{equation}
    \mathcal{G} = \mathcal{B}_0 \times \mathcal{B}^d,
    \qquad |\mathcal{G}| = G = M^{d+1}.
\end{equation}
It is stored as a matrix $\mathbf{B} \in \mathbb{R}^{G \times (d+1)}$
where row $g$ is the coefficient vector $\bsbeta^{(g)}$.

\section{Vectorized Welfare Evaluation}

\subsection{Single Auction Draw}

For a single auction draw $A^{(j)}$ with $n$ advertisers, define
\begin{equation}
    \mathbf{v} = \bigl(v_1,\ldots,v_n\bigr)^{\!\top} \in \mathbb{R}^n, \qquad
    \mathbf{e} = \bigl(e_1,\ldots,e_n\bigr)^{\!\top} \in \mathbb{R}^n,
\end{equation}
(dropping the superscript $j$ for clarity and writing $v, e$ for
$\tilde{v}, \tilde{e}$ within this section).
Construct the Vandermonde-style matrix of externality powers:
\begin{equation}
    \mathbf{E} \in \mathbb{R}^{n \times (d+1)}, \qquad
    E_{ip} = e_i^{\,p}, \quad i = 1,\ldots,n,\; p = 0,\ldots,d.
\end{equation}

\paragraph{Step 1: Evaluate $\tau$ for all grid points and all advertisers.}
\begin{equation}
    \mathbf{T} = \mathbf{B}\,\mathbf{E}^{\!\top} \in \mathbb{R}^{G \times n},
    \qquad T_{gi} = \tau(e_i;\,\bsbeta^{(g)}).
    \label{eq:tau_matrix}
\end{equation}
This is a single dense matrix multiply that computes $\tau(e_i)$ for every
grid point $g$ and every advertiser $i$ simultaneously.

\paragraph{Step 2: Determine admission.}
\begin{equation}
    \mathbf{A} \in \{0,1\}^{G \times n}, \qquad
    A_{gi} = \mathbf{1}[v_i \geq T_{gi}].
    \label{eq:admission}
\end{equation}

\paragraph{Step 3: Select top-$k$ admitted winners.}
Sort the $n$ advertisers by value descending:
let $\sigma$ be the permutation with $v_{\sigma(1)} \geq \cdots \geq v_{\sigma(n)}$.
Define the reordered matrices:
\begin{equation}
    \mathbf{v}^{\sigma} = (v_{\sigma(1)},\ldots,v_{\sigma(n)})^{\!\top},\quad
    \mathbf{e}^{\sigma} = (e_{\sigma(1)},\ldots,e_{\sigma(n)})^{\!\top},\quad
    \mathbf{A}^{\sigma}_{g,j} = A_{g,\sigma(j)}.
\end{equation}
Because advertisers are now in descending-$v$ order, the top-$k$ admitted
winners under grid point $g$ are exactly the first $k$ positions $j$ for
which $\mathbf{A}^{\sigma}_{g,j} = 1$.
This is identified via the cumulative admission count:
\begin{equation}
    C_{g,j} = \sum_{j'=1}^{j} A^{\sigma}_{g,j'}, \qquad
    \text{(number of admitted advertisers in positions } 1, \ldots, j\text{)}.
\end{equation}
The winner mask is:
\begin{equation}
    \mathbf{M} \in \{0,1\}^{G \times n}, \qquad
    M_{g,j} = \mathbf{1}[C_{g,j} \leq k]\cdot A^{\sigma}_{g,j}.
    \label{eq:winner_mask}
\end{equation}
Position $j$ is a winner under grid point $g$ if and only if it is admitted
\emph{and} is among the first $k$ admitted positions in the sorted order.

\paragraph{Step 4: Compute welfare.}
\begin{equation}
    \mathbf{w}^{(j)} = \mathbf{M}\,\mathbf{v}^{\sigma}
                     + \lambda\,\bigl(\mathbf{M}\,\mathbf{e}^{\sigma}\bigr)
                     \in \mathbb{R}^{G},
    \label{eq:auction_welfare}
\end{equation}
where $\lambda = R_e / R_v$ is the externality scale factor from
\eqref{eq:scaled_welfare} and the matrix--vector products
$\mathbf{M}\mathbf{v}^{\sigma}$ and $\mathbf{M}\mathbf{e}^{\sigma}$
sum the value and externality of the winners for each of the $G$ grid points.

\subsection{Averaging over Auction Draws}

Steps 1--4 are repeated for each of the $K$ auction draws.
The mean collateralized welfare over all draws is:
\begin{equation}
    \bar{\mathbf{w}} = \frac{1}{K} \sum_{j=1}^{K} \mathbf{w}^{(j)} \in \mathbb{R}^{G}.
    \label{eq:mean_welfare}
\end{equation}
In practice, $\bar{\mathbf{w}}$ is accumulated by adding $\mathbf{w}^{(j)}$
to a running sum and dividing by $K$ at the end, requiring only $O(G)$
additional memory regardless of $K$.

\section{Best Coefficient Selection}

The optimal grid point is:
\begin{equation}
    g^* = \operatorname*{argmax}_{g \in \{1,\ldots,G\}} \bar{w}_g,
    \qquad
    \bsbeta^* = \bsbeta^{(g^*)} \in \mathbb{R}^{d+1}.
    \label{eq:argmax}
\end{equation}
The coefficients $\bsbeta^*$ are the optimal tau coefficients in
normalized $(\te, \tilde{v})$ space.

\section{Coefficient Recovery in Original Space}

The stored result and all plots use the tau function in original $(e, v)$ space.
Since the normalization \eqref{eq:normalization} is an affine transform, the
polynomial $\tau(\te;\bsbeta^*)$ in normalized space corresponds to a polynomial
$\tau^*(e;\boldsymbol{\alpha}^*)$ in original space, found by substituting
$\te = \alpha_e e + \beta_e$ and $\tilde{v} = \alpha_v v + \beta_v$ where
\begin{equation}
    \alpha_e = \frac{2}{R_e}, \quad \beta_e = -1 - \alpha_e e_{\min}, \qquad
    \alpha_v = \frac{2}{R_v}, \quad \beta_v = -1 - \alpha_v v_{\min}.
\end{equation}
Given $\tilde{v} = \tau(\te;\bsbeta^*)$, we solve for $v$:
\begin{equation}
    v = \frac{\tau(\alpha_e e + \beta_e;\;\bsbeta^*) - \beta_v}{\alpha_v}
      = \tau^*(e;\,\boldsymbol{\alpha}^*).
    \label{eq:recovery}
\end{equation}
The polynomial composition $p(\alpha_e x + \beta_e)$ is computed exactly
via Horner's rule, and the affine scaling by $1/\alpha_v$ and shift by
$-\beta_v/\alpha_v$ produce $\boldsymbol{\alpha}^*$ as a plain vector of
coefficients in increasing-power order.
The result satisfies:
for all $(e, v)$ in original space,
$v \geq \tau^*(e;\boldsymbol{\alpha}^*)$ if and only if
$\tv \geq \tau(\te;\bsbeta^*)$.

\section{Algorithm Summary}

\begin{algorithm}[H]
\caption{Grid Search for Optimal Tau}
\noindent\textbf{Require:} Dataset $\{(v_t^{\text{score}}, e_t^{\text{score}})\}$, cost parameters
$\zeta_v, \zeta_e$, degree $d$, grid size $M$, slots $k$, draws $K$,
items per draw $n$, random seed
\begin{description}
\setlength{\itemsep}{1pt}\setlength{\parsep}{0pt}
\item[\normalfont 1:]  $v_t \gets \zeta_v v_t^{\text{score}}$;\quad $e_t \gets \zeta_e e_t^{\text{score}}$
\item[\normalfont 2:]  Fit independent scaler; compute $\tilde{v}_t, \tilde{e}_t$ via \eqref{eq:normalization}
\item[\normalfont 3:]  $\lambda \gets R_e / R_v$
\item[\normalfont 4:]  Sample $K$ draws $\mathcal{A} = \{A^{(j)}\}$ of $n$ normalized pairs each
\item[\normalfont 5:]  Build grid $\mathbf{B} \in \mathbb{R}^{G \times (d+1)}$,\quad $G = M^{d+1}$
\item[\normalfont 6:]  $\bar{\mathbf{w}} \gets \mathbf{0}_G$
\item[\normalfont 7:]  \textbf{for} $j = 1$ \textbf{to} $K$ \textbf{do}
\item[\normalfont 8:]  \hspace{1.5em}$\mathbf{E} \gets$ Vandermonde matrix of $e$-powers for $A^{(j)}$ \hfill \eqref{eq:tau_matrix}
\item[\normalfont 9:]  \hspace{1.5em}$\mathbf{T} \gets \mathbf{B}\,\mathbf{E}^\top$ \hfill (evaluate $\tau$ for all $g$, $i$)
\item[\normalfont 10:] \hspace{1.5em}$\mathbf{A} \gets [v_i \geq T_{gi}]$ \hfill \eqref{eq:admission}
\item[\normalfont 11:] \hspace{1.5em}Sort advertisers by $v$ descending; compute winner mask $\mathbf{M}$ \hfill \eqref{eq:winner_mask}
\item[\normalfont 12:] \hspace{1.5em}$\bar{\mathbf{w}} \gets \bar{\mathbf{w}} + \mathbf{M}\mathbf{v}^\sigma + \lambda\,\mathbf{M}\mathbf{e}^\sigma$ \hfill \eqref{eq:auction_welfare}
\item[\normalfont 13:] \textbf{end for}
\item[\normalfont 14:] $\bar{\mathbf{w}} \gets \bar{\mathbf{w}} / K$
\item[\normalfont 15:] $g^* \gets \operatorname{argmax}_g\, \bar{w}_g$;\quad $\bsbeta^* \gets \mathbf{B}_{g^*}$ \hfill \eqref{eq:argmax}
\item[\normalfont 16:] $\boldsymbol{\alpha}^* \gets \texttt{poly\_from\_normalized}(\bsbeta^*, e, v)$ \hfill \eqref{eq:recovery}
\item[\normalfont 17:] \textbf{return} $\boldsymbol{\alpha}^*$
\end{description}
\end{algorithm}

\section{Complexity}

The dominant cost per auction draw is the matrix multiply in step~8,
which costs $O(G \cdot n \cdot (d+1))$ floating-point operations.
The total cost over all draws is $O(K \cdot G \cdot n \cdot (d+1))$.
For the default parameters ($K = 500$, $n = 20$, $d = 1$, $M = 20$,
$G = 400$) this is approximately $8 \times 10^6$ operations,
well within the range of a single-threaded NumPy call.
For $d = 3$, $M = 20$ ($G = 160{,}000$) the count rises to
$\sim\!3.2 \times 10^9$, at which point the vectorized matrix
operations benefit noticeably from multi-threaded BLAS.
The memory requirement per draw is $O(G \cdot n)$ for
$\mathbf{T}$, $\mathbf{A}$, and $\mathbf{M}$, plus $O(G)$
for the running welfare sum $\bar{\mathbf{w}}$.

\end{document}
